\paper{非数学类}

\begin{problems}
\item[一、]（本题 42 分，共 6 小题，每小题 7 分）填空题：
\begin{enumerate}
  \item 已知可导函数 $f(x)$ 满足 $\displaystyle f(x)\cos x+2\int_0^x f(t)\sin t\dd t=x+1$，则 $f(x)=\underline{\hspace{3.5em}}$．
  \item 极限 $\displaystyle \lim_{n\to\infty}\sin^2\left(\pi\sqrt{n^2+n}\right)=\underline{\hspace{3.5em}}$．
  \item 设 $w=f(u,v)$ 具有二阶连续偏导数，且 $u=x-cy$，$v=x+cy$，其中 $c$ 为非零常数，则 $\displaystyle w_{xx}-\frac{1}{c^2}w_{yy}=\underline{\hspace{3.5em}}$．
  \item 设 $f(x)$ 具有二阶连续导数，且 $f(0)=f'(0)=0$，$f''(0)=6$，则 $\displaystyle \lim_{x\to 0}\frac{f(\sin^2 x)}{x^4}=\underline{\hspace{3.5em}}$．
  \item 不定积分 $\displaystyle I=\int\frac{\ee^{-\sin x}\sin 2x}{(1-\sin x)^2}\dd x=\underline{\hspace{3.5em}}$．
  \item 记曲面 $z^2=x^2+y^2$ 和 $z=\sqrt{4-x^2-y^2}$ 围成的空间区域为 $V$，则三重积分
  \[
    \iiint_V z\dd x\dd y\dd z=\underline{\hspace{10em}}.
  \]
\end{enumerate}

\item[二、]（本题 14 分）设二元函数 $f(x,y)$ 在平面上有连续的二阶偏导数，对任意角度 $\alpha$，定义一元函数 $g_\alpha(t)=f(t\cos\alpha,t\sin\alpha)$，若对任何 $\alpha$ 都有 $\displaystyle \frac{\dd g_\alpha(0)}{\dd t}=0$ 且 $\displaystyle \frac{\dd^2 g_\alpha(0)}{\dd t^2}>0$，证明：$f(0,0)$ 是 $f(x,y)$ 的极小值．

\item[三、]（本题 14 分）设曲线 $\Gamma$ 为曲线 $x^2+y^2+z^2=1$，$x+z=1$，$x\ge 0$，$y\ge 0$，$z\ge 0$ 上从点 $A(1,0,0)$ 到点 $B(0,0,1)$ 的一段．求曲线积分
\[
  I=\int_\Gamma y\dd x+z\dd y+x\dd z.
\]

\item[四、]（本题 15 分）设函数 $f(x)>0$ 且在实轴上连续，若对任意实数 $t$，有
\[
  \int_{-\infty}^{+\infty}\ee^{-|t-x|}f(x)\dd x\le 1.
\]
证明：$\forall a,b$，$a<b$，有 $\displaystyle \int_a^b f(x)\dd x\le\frac{b-a+2}{2}$．

\item[五、]（本题 15 分）设 $\{a_n\}$ 为一个数列，$p$ 为固定的正整数，若 $\lim\limits_{n\to\infty}(a_{n+p}-a_n)=\lambda$．证明：
\[
  \lim_{n\to\infty}\frac{a_n}{n}=\frac{\lambda}{p}.
\]
\end{problems}
